% Translated from sections/05-results.tex on 2026-09-23. The English file is
% canonical. Change it first, then carry the change here.
%
% Findings only, as in the English file. Interpretation is section 6.
\section{Résultats}
\label{sec:results}

Chaque partie s'ouvre sur son affirmation, son unité d'analyse, son
système de comparaison et sa méthode d'incertitude, puis rapporte les
résultats. E1, E3 et E4 n'y figurent pas. Le jeu de référence qu'elles
exigent existe désormais (\goldPages{}~pages, section~\ref{sec:protocol}), mais
aucune des trois n'a été exécutée.

\subsection{E5: localiser un extrait du \emph{Trésor}}
\label{sec:e5}

\textbf{Affirmation.} Un extrait du \emph{Trésor} peut être localisé dans le
texte de base. \textbf{Unité.} La pièce, au sein d'une cohorte.
\textbf{Comparaison.} Aucune n'est publiée. Le critère préenregistré en tient
lieu. \textbf{Incertitude.} Bootstrap sur les pièces.

\begin{table}[t]
  \centering
  % Generated by src/amadis_htr/report_macros.py from eval/results/*.csv.
% Do not edit: rerun eval/run_all.sh instead.
%
% Empty until the evaluation that feeds it runs. Citing a figure before it
% is measured is an undefined control sequence, which stops the build.
% That is the point.
\begin{tabular}{lrrrr}
\toprule
Lot & Extraits & Bon Livre & Chapitre vérifiable & Bon chapitre \\
\midrule
lot 1 & 317 & 315 (99,4\,\%) & 125 & 96 (76,8\,\%) \\
lot 2 & 101 & 100 (99,0\,\%) & 52 & 38 (73,1\,\%) \\
\bottomrule
\end{tabular}

  \caption{E5, attributions certaines seulement. Les deux cohortes ne sont
    jamais regroupées, car le critère a été préenregistré sur le lot du
    classeur quand il était le seul à exister.}
  \label{tab:e5}
\end{table}

Il y a trois dénominateurs, et ils diffèrent. Les pièces de la cohorte, les
pièces certaines pour un chiffre au niveau du Livre, et les pièces dont le
numéro de chapitre imprimé est récupérable pour un chiffre au niveau du
chapitre. Ces dernières sont \locChapterScoreableWorkbook{} sur
\locPiecesWorkbook{} dans la cohorte soumise au critère. Parmi les autres,
\locChapterUnavailableWorkbook{}~pièces sont tombées sur un intitulé imprimé
dont aucun numéro de chapitre n'a pu être tiré, et
\locChapterUnreferencedWorkbook{} a été cataloguée dans un Livre sans chapitre
nommé.

Le tableau~\ref{tab:e5} donne le résultat. L'exactitude au niveau du Livre est
de \locLivreAccuracyWorkbook{} sur la cohorte classeur et de
\locLivreAccuracyCatalogue{} sur la cohorte catalogue. L'exactitude au niveau
du chapitre est de \locChapterAccuracyWorkbook{} et de
\locChapterAccuracyCatalogue{}. Le critère, écrit avant l'évaluation de
l'apparieur, se lit ainsi: «\textenglish{\emph{Correct LIVRE well above 90\%
means the all-24-Livres search holds \dots{} Proceed to phase B}}».\footnote{Un
taux de Livre exact nettement supérieur à 90\,\% signifie que la recherche sur
les 24~Livres tient \dots{} Passer à la phase~B.}

\begin{figure}[t]
  \centering
  \includegraphics[width=\linewidth]{generated/fr/figures/e5-sweep.pdf}
  \caption{E5, seuil d'acceptation contre précision (en haut) et couverture
    (en bas). Les seuils qui n'acceptent rien sont omis, la précision sur un
    ensemble accepté vide n'étant pas définie.}
  \label{fig:e5-sweep}
\end{figure}

La figure~\ref{fig:e5-sweep} fait varier le seuil d'acceptation. La précision
reste proche de 100\,\% à tous les seuils dans les deux cohortes. La courbe ne
désigne donc pas de point de fonctionnement, et c'est la couverture qui le
fait. Elle reste supérieure à 98\,\% jusqu'à 0,86 dans la cohorte classeur et
jusqu'à 0,85 dans la cohorte catalogue, puis décroche. À 0,90 elle n'est plus
que de 73\,\% et 53\,\%. L'effet le plus important de E5 est l'écart entre
l'exactitude au niveau du Livre et celle au niveau du chapitre.

\subsection{E2: la règle d'acceptation par ancrage}
\label{sec:e2}

\textbf{Affirmation.} La règle d'acceptation empêche une passe générative de
réécrire le texte. \textbf{Unité.} La ligne soumise, regroupée par page.
\textbf{Comparaison.} La correction désactivée, qui n'est autre que la sortie
de reconnaissance figée. Les \ocrBooks{}~livres ont tous enregistré zéro
mot corrigé. \textbf{Incertitude.} Bootstrap stratifié sur les pages,
10\,000~tirages, à 95\,\%.

Sur l'ensemble du corpus, la passe examinerait \corrReachLinesWithSpan{} des
\corrReachCorrectable{}~lignes corrigeables (\corrReachLineShare{}), portant
\corrReachSpans{}~segments suspects, sur \corrReachPagesWithSpan{} des
\corrReachPages{}~pages (\corrReachPageShare{}). Ce décompte est déterministe.
Il ne dépend que des confiances du reconnaisseur, sans aucun modèle.
L'échantillon compte \corrOffered{}~lignes sur \corrPages{}~pages, tirées à
raison de 200 par famille typographique avec une graine versionnée.

\begin{table}[t]
  \centering
  % Generated by src/amadis_htr/report_macros.py from eval/results/*.csv.
% Do not edit: rerun eval/run_all.sh instead.
%
% Empty until the evaluation that feeds it runs. Citing a figure before it
% is measured is an undefined control sequence, which stops the build.
% That is the point.
\begin{tabular}{lrrrrr}
\toprule
Cohorte & Lignes & Acceptées & Réécritures & Échecs de protocole & Repli par recherche \\
\midrule
famille A & 450 & 414 (92,0\,\%) & 16 & 20 & 325 \\
famille B & 500 & 478 (95,6\,\%) & 10 & 12 & 324 \\
échantillon & 950 & 892 & 26 & 32 & 649 \\
\bottomrule
\end{tabular}

  \caption{E2 par famille typographique. \emph{Lignes} compte les lignes
    soumises à la passe, \emph{Réécritures} les
    réponses refusées pour avoir modifié le texte hors des marqueurs. La ligne
    de l'échantillon ne porte que des effectifs. Le tirage est égal par
    famille alors que la base de sondage ne l'est pas, si bien qu'un taux
    calculé sur cette ligne donnerait à la famille~A plus du double de son
    poids. Le taux du corpus est pondéré par la base de sondage.}
  \label{tab:e2}
\end{table}

Sur les \corrOffered{}~lignes soumises, \corrAccepted{} ont été acceptées. Pondéré
aux lignes suspectes du corpus, le taux d'acceptation est de \corrAcceptRate{}
(IC à 95\,\% de \corrAcceptLow{} à \corrAcceptHigh{}). Parmi les refus,
\corrRewritesCaught{} tenaient à une modification du texte hors des marqueurs.
Cela représente \corrRewriteRate{} des lignes soumises après pondération, sur
\corrRewriteSegments{}~segments et \corrRewriteEdits{}~caractères, dont
\corrRewritesSubstantial{} déplaçaient plus de deux caractères. Les
\corrProtocolFailures{} autres refus sont des réponses qui ne sont jamais
devenues des modifications candidates, et ils ne pèsent en rien sur
l'affirmation. La figure~\ref{fig:e2-confidence} compare la confiance du
reconnaisseur sur les lignes acceptées et refusées.

\begin{figure}[t]
  \centering
  \includegraphics[width=\linewidth]{generated/fr/figures/e2-confidence.pdf}
  \caption{Confiance minimale d'un caractère sur une ligne, lignes acceptées
    contre lignes refusées.}
  \label{fig:e2-confidence}
\end{figure}

Parmi les lignes soumises, \corrOnFindFallback{} (\corrFindFallbackShare{}) ont
emprunté le repli de localisation des segments décrit en
section~\ref{sec:errors}. Sur ces lignes, une acceptation ne prouve pas que la
règle a tenu.

L'exécution a utilisé quatre processus concurrents. Le protocole enregistré
supposait que cela changerait le débit et non le verdict. Un rejeu de
\corrReplayLines{}~lignes tirées d'un échantillon réduit de
\corrReplayPages{}~pages, avec graine et un seul processus, a reproduit
\corrReplayAgree{} des \corrReplayLines{}~verdicts. C'est une donnée là où il
n'y avait qu'une hypothèse, et non une preuve que le décodage est
déterministe.

L'acceptation n'est pas un taux de modification. Parmi les lignes acceptées,
\corrAcceptedUnchanged{} sont revenues avec tous leurs segments intacts, et
\corrChangedSpans{}~segments ont réellement été modifiés. \textbf{Ce n'est pas
non plus une exactitude.} Comparer la correction activée à la correction
désactivée exige une référence pour les pages tirées par E2, et le jeu de
référence n'en couvre aucune. E2 rapporte ce qu'a fait la règle, et ne dit
rien de l'effet des modifications retenues sur la qualité du texte.

\subsection{E6: débit et structure}
\label{sec:e6}

\textbf{Affirmation.} Descriptive seulement. \textbf{Unité.} Le livre et la
page. \textbf{Comparaison.} Aucune. \textbf{Incertitude.} Aucune, car ce sont
des décomptes.

\begin{table}[t]
  \centering
  % Generated by src/amadis_htr/report_macros.py from eval/results/*.csv.
% Do not edit: rerun eval/run_all.sh instead.
%
% Empty until the evaluation that feeds it runs. Citing a figure before it
% is measured is an undefined control sequence, which stops the build.
% That is the point.
\begin{tabular}{lrrrrr}
\toprule
Cohorte & Livres & Pages & Échecs & Heures & s/page \\
\midrule
famille A & 12 & 3\,658 & 0 & 1,53 & 1,509 \\
famille B & 12 & 10\,453 & 0 & 4,11 & 1,417 \\
corpus & 24 & 14\,111 & 0 & 5,65 & 1,441 \\
\bottomrule
\end{tabular}

  \caption{Débit de E6. Un même modèle, un même matériel et une même largeur
    de travail pour les \ocrBooks{}~livres, ce qui rend le regroupement
    légitime.}
  \label{tab:e6}
\end{table}

Les \ocrBooks{}~livres, soit \ocrPages{}~pages, ont été traités avec
\ocrPagesFailed{}~échec en \ocrHours{}~heures, à \ocrSecondsPerPage{}~s par page. La
figure~\ref{fig:e6-throughput} montre la dispersion par livre. Les familles
diffèrent peu, avec \ocrSecondsPerPageFamilyA{}~s contre
\ocrSecondsPerPageFamilyB{}~s.

\begin{figure}[t]
  \centering
  \includegraphics[width=\linewidth]{generated/fr/figures/e6-throughput.pdf}
  \caption{Secondes par page, par livre.}
  \label{fig:e6-throughput}
\end{figure}

\textbf{La structure ne porte que sur la famille~A, et elle a été vérifiée à la
main.} Les livres~13 à~24 n'impriment pas de table des matières. Leur
calibration ne comporte donc pas de bloc de correspondance, et il n'existe pas
de chiffre structurel pour tout le corpus. Sur la famille~A, l'exécution
retrouve \structMatchedA{} des \structChaptersA{}~chapitres imprimés
(\structRecallHandCheckedA{}) avec une précision de
\structPrecisionHandCheckedA{}. Elle en manque \structMissingA{} et propose
\structSpuriousA{}~débuts de chapitre que la table imprimée ne porte pas. Ces
deux taux ne décrivent pas le détecteur seul. Le fichier de corrections de
chaque livre contient des décisions de chapitre prises à la main par un
lecteur, et la calibration a été exécutée après elles. Ce qui est rapporté est
donc le détecteur corrigé. Les décomptes qui les accompagnent sont exactement
ce qu'ils disent être.
